%! AP Statistics — Unit 2, Lesson 2.4 — Student Packet Cover Sheet
\documentclass[10pt]{article}
\usepackage{apstats-article}
\usepackage{apstats-boxes}

%=============================================================================
\begin{document}

% ── Full-bleed navy banner ────────────────────────────────────────────────────
\begin{tikzpicture}[remember picture, overlay]
  \fill[navy] (current page.north west)
    rectangle ([yshift=-0.9in]current page.north east);
\end{tikzpicture}
\vspace{-0.8in}

\begin{center}
  {\color{white}\textbf{\LARGE AP Statistics}}\\[4pt]
  {\color{skymid}\large\textbf{Unit 2: Exploring Two-Variable Data}}\\[6pt]
  {\color{white}\textbf{\large Lesson 2.4 \quad Representing the Relationship Between Two Quantitative Variables}}
\end{center}

\vspace{0.22in}

\namedateperiod

\vspace{0.18in}

% ── LEARNING TARGETS ─────────────────────────────────────────────────────────
\begin{learningtargetbox}
\small
\begin{enumerate}[leftmargin=1.5em, itemsep=5pt]
  \item \ldots{} identify the explanatory and response variables and place them on the correct
        axes of a scatterplot. \hfill\textcolor{slate}{\small(UNC-1.S)}
  \item \ldots{} construct a correctly labeled scatterplot from a two-variable dataset.
        \hfill\textcolor{slate}{\small(UNC-1.S)}
  \item \ldots{} describe the \textbf{direction}, \textbf{form}, and \textbf{strength} of the
        association shown in a scatterplot, and identify any \textbf{outliers}.
        \hfill\textcolor{slate}{\small(UNC-1.T)}
  \item \ldots{} write a complete DOFS description in context.
        \hfill\textcolor{slate}{\small(UNC-1.T)}
\end{enumerate}
\end{learningtargetbox}

\vspace{0.14in}

% ── TABLE OF CONTENTS ────────────────────────────────────────────────────────
\begin{tocbox}
\small
\renewcommand{\arraystretch}{1.5}
\begin{tabularx}{\linewidth}{@{} c l X r @{}}
\rowcolor{sky}
\textbf{\#} & \textbf{Component} & \textbf{Description} & \textbf{Score} \\
\midrule
1 & Warm-Up       & Spiral review --- explanatory \& response variables              & \blank{1.2cm} \\
2 & Guided Notes  & Scatterplots and the DOFS framework                              & \blank{1.2cm} \\
3 & Group Activity & Construct and describe scatterplots from real-world data        & \blank{1.2cm} \\
4 & Exit Ticket   & Write a complete DOFS description from a given scatterplot       & \blank{1.2cm} \\
5 & Homework      & Multi-context scatterplot construction and description practice   & \blank{1.2cm} \\
\midrule
  & \multicolumn{2}{r}{\textbf{Total}} & \blank{1.2cm} \\
\end{tabularx}
\end{tocbox}

\vspace{0.14in}

% ── KEEP IN MIND ─────────────────────────────────────────────────────────────
\begin{remindbox}
\small
A scatterplot shows the relationship between two quantitative variables. Here are the key ideas.

\vspace{6pt}
\begin{tabularx}{\linewidth}{@{} >{\centering\arraybackslash\columncolor{goldbg}}p{0.06\linewidth}
                                    X @{}}
\textbf{\large!} &
\textbf{Explanatory variable goes on the $x$-axis.}\enspace
The explanatory variable may help explain or predict the response variable ($y$-axis).
Axis placement matters for regression in later lessons. \\[6pt]
\textbf{\large!} &
\textbf{Describe using DOFS.}\enspace
\textbf{D}irection (positive/negative/none), \textbf{F}orm (linear/curved/none),
\textbf{S}trength (strong/moderate/weak), \textbf{O}utliers (any unusual points). \\[6pt]
\textbf{\large!} &
\textbf{Always write in context.}\enspace
Name the variables and the individuals. ``There is a positive association'' is incomplete ---
say \emph{what} is positively associated and \emph{for whom}. \\[6pt]
\textbf{\large!} &
\textbf{Positive does not mean good.}\enspace
Positive association simply means that as $x$ increases, $y$ tends to increase.
It has nothing to do with whether the relationship is desirable.
\end{tabularx}
\end{remindbox}

\end{document}
